\documentclass{article}
\usepackage{amsmath,amssymb,amsthm}
\usepackage{algorithm}
\usepackage{algpseudocode}
\usepackage{bm}
\usepackage{enumitem}
\usepackage{hyperref}
\newtheorem{theorem}{Theorem}
\newtheorem{lemma}{Lemma}
\newtheorem{assumption}{Assumption}
\newcommand{\EE}{\mathbb{E}}
\newcommand{\RR}{\mathbb{R}}
\newcommand{\norm}[1]{\left\|#1\right\|}
\newcommand{\diag}{\operatorname{diag}}
\begin{document}

\section*{AdamW (Adam with Decoupled Weight Decay)}

\section*{Mathematical Formulation}
Let $f:\RR^{d}\!\to\!\RR$ be the (possibly stochastic) objective and let
$\{\xi_t\}_{t\ge 1}$ be a sequence of i.i.d.\ data samples.
Denote the (stochastic) gradient at iteration $t$ by
\[
g_t \;:=\; \nabla f(w_t;\xi_t)\in\RR^{d}.
\]

Hyper‑parameters:
\begin{itemize}[noitemsep]
    \item learning rate $\eta>0$,
    \item exponential decay rates $\beta_1,\beta_2\in[0,1)$,
    \item small constant $\epsilon>0$,
    \item weight‑decay coefficient $\lambda\ge 0$.
\end{itemize}

The algorithm maintains two moving averages:
\[
\begin{aligned}
m_t &= \beta_1 m_{t-1} + (1-\beta_1) g_t, \qquad m_0 = 0,\\[4pt]
v_t &= \beta_2 v_{t-1} + (1-\beta_2) g_t\odot g_t,\qquad v_0 = 0,
\end{aligned}
\]
where $\odot$ denotes element‑wise product.

Bias‑corrected moments:
\[
\hat m_t \;=\; \frac{m_t}{1-\beta_1^{\,t}},\qquad
\hat v_t \;=\; \frac{v_t}{1-\beta_2^{\,t}}.
\]

The **decoupled weight‑decay** step is applied directly to the parameters,
followed by the Adam step:
\[
\boxed{
\begin{aligned}
w_{t+\frac12}&=w_t-\eta\lambda w_t,\\[2pt]
w_{t+1}&= w_{t+\frac12}
          -\eta\,\frac{\hat m_t}{\sqrt{\hat v_t}+\epsilon},
\end{aligned}}
\]
where the square‑root and division are taken element‑wise.

\section*{Pseudocode}
\begin{algorithm}[H]
\caption{AdamW}
\begin{algorithmic}[1]
\State \textbf{Input:} initial point $w_0\in\RR^{d}$, learning rate $\eta$, 
      decay rates $\beta_1,\beta_2$, weight‑decay $\lambda$, $\epsilon>0$
\State Initialize $m_0\gets 0$, $v_0\gets 0$
\For{$t=1,2,\dots,T$}
    \State Sample $\xi_t$ and compute stochastic gradient $g_t\gets\nabla f(w_t;\xi_t)$
    \State $m_t \gets \beta_1 m_{t-1} + (1-\beta_1) g_t$
    \State $v_t \gets \beta_2 v_{t-1} + (1-\beta_2) (g_t\odot g_t)$
    \State $\hat m_t \gets m_t/(1-\beta_1^{\,t})$
    \State $\hat v_t \gets v_t/(1-\beta_2^{\,t})$
    \State $w_{t+\frac12}\gets w_t - \eta\lambda w_t$
    \State $w_{t+1}\gets w_{t+\frac12}
            - \eta\,\dfrac{\hat m_t}{\sqrt{\hat v_t}+\epsilon}$
\EndFor
\State \textbf{Output:} $w_{T+1}$
\end{algorithmic}
\end{algorithm}

\section*{Dry Run Example}
Consider the one‑dimensional quadratic 
\[
f(w)=(w-3)^2,\qquad \nabla f(w)=2(w-3).
\]
Choose hyper‑parameters
\[
\eta=0.1,\;\beta_1=0.9,\;\beta_2=0.999,\;
\lambda=0.01,\;\epsilon=10^{-8},
\]
and initialise $w_0=0$, $m_0=v_0=0$.

\begin{center}
\begin{tabular}{c|c|c|c|c|c}
$t$ & $g_t$ & $m_t$ & $v_t$ & $w_{t+1}$ & $f(w_{t+1})$\\\hline
1 & $2(0-3)=-6$ &
  $0.1\cdot(-6)=-0.6$ &
  $0.001\cdot36=0.036$ &
  $w_{1}=0-0.1\cdot0.01\cdot0
          -0.1\frac{-0.6}{\sqrt{0.036}+10^{-8}}
          =0+0.1\frac{0.6}{0.189736}=0.316$ &
  $(0.316-3)^2\approx7.20$\\\hline
2 & $2(0.316-3)=-5.368$ &
  $0.9(-0.6)+(0.1)(-5.368)=-1.0868$ &
  $0.999(0.036)+(0.001)(28.81)=0.06481$ &
  $w_{2}=0.316-0.1\cdot0.01\cdot0.316
          -0.1\frac{-1.0868}{\sqrt{0.06481}+10^{-8}}
          \approx0.316-0.000316+0.1\frac{1.0868}{0.2546}
          =0.727$ &
  $(0.727-3)^2\approx5.18$\\\hline
3 & $2(0.727-3)=-4.546$ &
  $0.9(-1.0868)+(0.1)(-4.546)=-1.5142$ &
  $0.999(0.06481)+(0.001)(20.66)=0.08546$ &
  $w_{3}=0.727-0.1\cdot0.01\cdot0.727
          -0.1\frac{-1.5142}{\sqrt{0.08546}+10^{-8}}
          \approx0.727-0.000727+0.1\frac{1.5142}{0.2923}
          =1.087$ &
  $(1.087-3)^2\approx3.68$
\end{tabular}
\end{center}
The iterates move towards the minimiser $w^\star=3$, and the objective value decreases.

\newpage
\section*{Assumptions}
\begin{assumption}[Smoothness]\label{as:smooth}
$f$ has $L$‑Lipschitz continuous gradients, i.e.
\[
\forall w,w'\in\RR^{d}:\qquad
\|\nabla f(w)-\nabla f(w')\|\le L\|w-w'\|.
\]
\end{assumption}
\begin{assumption}[Bounded Stochastic Gradient Moments]\label{as:bounded}
There exist constants $G,\sigma>0$ such that for all $t$,
\[
\EE\!\left[\|g_t\|^2\mid \mathcal{F}_{t-1}\right]\le G^2,
\qquad
\EE\!\left[\|g_t-\nabla f(w_t)\|^2\mid \mathcal{F}_{t-1}\right]\le\sigma^2,
\]
where $\mathcal{F}_{t-1}$ is the $\sigma$‑algebra generated by
$\{w_0,\dots,w_{t-1},\xi_1,\dots,\xi_{t-1}\}$.
\end{assumption}
\begin{assumption}[Hyper‑parameter Range]\label{as:hyper}
The decay rates satisfy
\[
0<\beta_1<\sqrt{\beta_2}<1,\qquad
\eta\le \frac{1}{L}\sqrt{1-\beta_2},\qquad
\lambda\ge0,\qquad\epsilon>0.
\]
\end{assumption}

\section*{Main Convergence Theorem}
\begin{theorem}[Non‑convex convergence of AdamW]\label{thm:adamw}
Under Assumptions~\ref{as:smooth}–\ref{as:hyper},
let $\{w_t\}_{t\ge0}$ be the sequence generated by AdamW with a
non‑increasing learning‑rate schedule
\[
\eta_t \;=\; \frac{\eta}{\sqrt{t}}\,,\qquad t\ge1 .
\]
Define the averaged squared gradient norm
\[
\Phi_T \;:=\; \frac{1}{T}\sum_{t=1}^{T}
        \EE\!\left[\|\nabla f(w_t)\|^2\right].
\]
Then for any $T\ge 1$,
\[
\Phi_T \;\le\;
\underbrace{\frac{2L\,\bigl(f(w_0)-f^\star\bigr)}{\eta\sqrt{T}}}_{\text{optimization term}}
\;+\;
\underbrace{
\frac{2\eta\,G^2\,(1+\lambda)}{(1-\beta_1)(1-\beta_2)^{1/2}}
\frac{\log T}{\sqrt{T}}
}_{\text{variance/adaptivity term}} .
\]
Consequently,
\[
\Phi_T = \mathcal{O}\!\bigl((\log T)/\sqrt{T}\bigr).
\]
\end{theorem}

\section*{Theoretical Properties}
\begin{enumerate}[noitemsep]
    \item \textbf{Stability.} The denominator $\sqrt{\hat v_t}+\epsilon$ prevents division by zero and bounds the effective step size between
    $\displaystyle\frac{\eta_t}{\sqrt{G^2+\epsilon}}$ and $\displaystyle\frac{\eta_t}{\epsilon}$.
    \item \textbf{Hyper‑parameter Sensitivity.}
    The condition $\beta_1<\sqrt{\beta_2}$ (used in Lemma~\ref{lem:moment}) guarantees that the bias‑corrected moments do not explode; it is the same regime in which standard Adam enjoys convergence guarantees.
    \item \textbf{Comparison to Baselines.}
    When $\lambda=0$ the bound reduces to the classical Adam rate
    $O((\log T)/\sqrt{T})$ (see \cite{Reddi2018}).  The extra term
    $2\eta\lambda G^2/(1-\beta_1)(1-\beta_2)^{1/2}$ reflects the
    explicit $L_2$ regularisation induced by weight decay.
\end{enumerate}

\section*{Discussion}
The theorem shows that, despite the adaptive per‑coordinate scaling and the decoupled weight decay, AdamW retains the same sub‑linear convergence rate as vanilla Adam for smooth non‑convex objectives. The logarithmic factor stems from the adaptivity of $v_t$; if $v_t$ were kept constant (i.e. SGD with weight decay) the rate would improve to $O(1/\sqrt{T})$.

\newpage
\section*{Preliminaries (Definitions and Lemmas)}
\begin{lemma}[Bound on bias‑corrected second moment]\label{lem:moment}
Under Assumptions~\ref{as:bounded} and \ref{as:hyper},
for any coordinate $i\in\{1,\dots,d\}$ and any $t\ge1$,
\[
\EE\!\left[\hat v_{t,i}\mid\mathcal{F}_{t-1}\right]
\;\le\;
\frac{G^2}{1-\beta_2}.
\]
\end{lemma}
\begin{proof}
From the definition $v_{t,i}= \beta_2 v_{t-1,i} + (1-\beta_2) g_{t,i}^2$,
taking expectation conditioned on $\mathcal{F}_{t-1}$ yields
\[
\EE[v_{t,i}\mid\mathcal{F}_{t-1}]
 = \beta_2 v_{t-1,i} + (1-\beta_2)\EE[g_{t,i}^2\mid\mathcal{F}_{t-1}]
 \le \beta_2 v_{t-1,i} + (1-\beta_2)G^2 .
\]
Iterating the recursion gives
\[
\EE[v_{t,i}\mid\mathcal{F}_{t-1}]
 \le G^2\bigl(1-\beta_2^{\,t}\bigr).
\]
Dividing by $1-\beta_2^{\,t}$ (bias correction) yields the claim. ∎
\end{proof}

\begin{lemma}[Descent Lemma for $L$‑smooth functions]\label{lem:descent}
If $f$ satisfies Assumption~\ref{as:smooth}, then for any $w,w'\in\RR^{d}$,
\[
f(w')\le f(w) + \langle\nabla f(w),w'-w\rangle
          +\frac{L}{2}\|w'-w\|^2 .
\]
\end{lemma}
\begin{proof}
Standard from $L$‑smoothness; see e.g. \cite{Nesterov2004}. ∎
\end{proof}

\section*{Main Proof (Step‑by‑Step Derivation)}
We prove Theorem~\ref{thm:adamw}.  All expectations are taken with respect to the randomness up to iteration $t$ unless otherwise noted.

\begin{enumerate}
\item[Step 1.] \textbf{Apply the descent lemma} (\ref{lem:descent}) with $w'=w_{t+1}$ and $w=w_t$:
\[
f(w_{t+1}) \le f(w_t) + \langle\nabla f(w_t),w_{t+1}-w_t\rangle
                + \frac{L}{2}\|w_{t+1}-w_t\|^2 .
\tag{1}
\]

\item[Step 2.] \textbf{Insert the AdamW update.}
Recall $w_{t+1}= w_t - \eta_t\lambda w_t - \eta_t\frac{\hat m_t}{\sqrt{\hat v_t}+\epsilon}$.
Hence
\[
w_{t+1}-w_t = -\eta_t\lambda w_t - \eta_t\frac{\hat m_t}{\sqrt{\hat v_t}+\epsilon}.
\tag{2}
\]

\item[Step 3.] \textbf{Expand the inner product in (1).}
Using (2),
\[
\begin{aligned}
\langle\nabla f(w_t),w_{t+1}-w_t\rangle
&= -\eta_t\lambda\langle\nabla f(w_t),w_t\rangle
   -\eta_t\Big\langle\nabla f(w_t),
      \frac{\hat m_t}{\sqrt{\hat v_t}+\epsilon}\Big\rangle .
\end{aligned}
\tag{3}
\]

\item[Step 4.] \textbf{Expand the squared norm term.}
From (2),
\[
\begin{aligned}
\|w_{t+1}-w_t\|^2
&= \eta_t^2\Big\|\lambda w_t +
   \frac{\hat m_t}{\sqrt{\hat v_t}+\epsilon}\Big\|^2\\
&\le 2\eta_t^2\Bigl(
   \lambda^2\|w_t\|^2
   +\Big\|\frac{\hat m_t}{\sqrt{\hat v_t}+\epsilon}\Big\|^2\Bigr)
\tag{4}
\end{aligned}
\]
where the inequality follows from $\|a+b\|^2\le2(\|a\|^2+\|b\|^2)$.

\item[Step 5.] \textbf{Bound the adaptive denominator.}
For any coordinate $i$,
\[
\frac{1}{\sqrt{\hat v_{t,i}}+\epsilon}
\;\le\;
\frac{1}{\epsilon},
\qquad
\frac{1}{\sqrt{\hat v_{t,i}}+\epsilon}
\;\ge\;
\frac{1}{\sqrt{G^2/(1-\beta_2)}+\epsilon}
\;=: \;c_{\min}>0 .
\tag{5}
\]

\item[Step 6.] \textbf{Upper bound the second term in (3).}
Using Cauchy–Schwarz and (5):
\[
\begin{aligned}
\Big\langle\nabla f(w_t),
      \frac{\hat m_t}{\sqrt{\hat v_t}+\epsilon}\Big\rangle
&\ge
c_{\min}\,\langle\nabla f(w_t),\hat m_t\rangle .
\end{aligned}
\tag{6}
\]

\item[Step 7.] \textbf{Relate $\hat m_t$ to the true gradient.}
Recall $m_t = \beta_1 m_{t-1}+(1-\beta_1)g_t$.
Taking expectation conditioned on $\mathcal{F}_{t-1}$,
\[
\EE[m_t\mid\mathcal{F}_{t-1}]
 = \beta_1 m_{t-1} + (1-\beta_1)\nabla f(w_t).
\]
Iterating yields
\[
\EE[m_t\mid\mathcal{F}_{t-1}]
 = (1-\beta_1^t)\nabla f(w_t) + \beta_1^t m_0 .
\]
Since $m_0=0$, bias‑correction gives
\[
\EE[\hat m_t\mid\mathcal{F}_{t-1}]
 = \nabla f(w_t).
\tag{7}
\]

\item[Step 8.] \textbf{Take total expectation of (1).}
Using (3)–(7) and noting that $\EE[\langle\nabla f(w_t),\hat m_t\rangle\mid\mathcal{F}_{t-1}]
 = \|\nabla f(w_t)\|^2$, we obtain
\[
\begin{aligned}
\EE\bigl[f(w_{t+1})\bigr]
&\le \EE\bigl[f(w_t)\bigr]
   -\eta_t\lambda\EE\bigl[\langle\nabla f(w_t),w_t\rangle\bigr]
   -\eta_t c_{\min}\EE\bigl[\|\nabla f(w_t)\|^2\bigr] \\
&\quad + L\eta_t^2\Bigl(
      \lambda^2\EE\bigl[\|w_t\|^2\bigr]
      + \EE\bigl[\big\|\tfrac{\hat m_t}{\sqrt{\hat v_t}+\epsilon}\big\|^2\bigr]\Bigr).
\end{aligned}
\tag{8}
\]

\item[Step 9.] \textbf{Bound the adaptive norm term.}
From Lemma~\ref{lem:moment} and $\|\hat m_t\|\le\|g_t\|/(1-\beta_1)$,
\[
\Big\|\frac{\hat m_t}{\sqrt{\hat v_t}+\epsilon}\Big\|^2
\le
\frac{\|g_t\|^2}{(1-\beta_1)^2\,c_{\min}^2}
\le
\frac{G^2}{(1-\beta_1)^2\,c_{\min}^2}.
\tag{9}
\]

\item[Step 10.] \textbf{Bound the weight‑decay inner product.}
Using Cauchy–Schwarz,
\[
|\langle\nabla f(w_t),w_t\rangle|
\le \|\nabla f(w_t)\|\;\|w_t\|
\le \tfrac12\bigl(\|\nabla f(w_t)\|^2+\|w_t\|^2\bigr).
\tag{10}
\]

\item[Step 11.] \textbf{Collect terms.}
Insert (9) and (10) into (8) and drop the non‑negative term $-\eta_t\lambda\,\tfrac12\|\nabla f(w_t)\|^2$ (which only improves the bound):
\[
\begin{aligned}
\EE\bigl[f(w_{t+1})\bigr]
&\le \EE\bigl[f(w_t)\bigr]
   -\underbrace{\Bigl(\eta_t c_{\min}
          -\tfrac12\eta_t\lambda\Bigr)}_{\displaystyle\alpha_t}
      \EE\bigl[\|\nabla f(w_t)\|^2\bigr] \\
&\quad + \underbrace{L\eta_t^2\lambda^2}_{\displaystyle\beta_t}
      \EE\bigl[\|w_t\|^2\bigr]
   + \underbrace{L\eta_t^2}_{\displaystyle\gamma_t}
      \frac{G^2}{(1-\beta_1)^2c_{\min}^2}.
\end{aligned}
\tag{11}
\]

\item[Step 12.] \textbf{Discard the $\EE[\|w_t\|^2]$ term.}
Since $\lambda\ge0$, the term $L\eta_t^2\lambda^2\EE[\|w_t\|^2]$
is non‑negative; dropping it yields a looser but simpler inequality:
\[
\EE\bigl[f(w_{t+1})\bigr]
\le \EE\bigl[f(w_t)\bigr]
   -\alpha_t\EE\bigl[\|\nabla f(w_t)\|^2\bigr]
   +\gamma_t\frac{G^2}{(1-\beta_1)^2c_{\min}^2}.
\tag{12}
\]

\item[Step 13.] \textbf{Choose the learning‑rate schedule.}
Set $\eta_t = \eta/\sqrt{t}$.
Because $c_{\min}$ depends only on $G,\beta_2,\epsilon$, we have
$\alpha_t = \eta_t\bigl(c_{\min}-\lambda/2\bigr)
            \ge \frac{\eta}{2\sqrt{t}}c_{\min}$,
provided $\lambda\le c_{\min}$ (which holds for the usual
small weight‑decay values).  Define
\[
C_1:=\frac{c_{\min}}{2},\qquad
C_2:=\frac{G^2}{(1-\beta_1)^2c_{\min}^2}.
\]

Thus (12) becomes
\[
\EE\bigl[f(w_{t+1})\bigr]
\le \EE\bigl[f(w_t)\bigr]
   -\frac{\eta C_1}{\sqrt{t}}\,
     \EE\bigl[\|\nabla f(w_t)\|^2\bigr]
   +\frac{\eta^2 L C_2}{t}.
\tag{13}
\]

\item[Step 14.] \textbf{Sum (13) from $t=1$ to $T$.}
Telescoping the left‑hand side yields
\[
\EE\bigl[f(w_{T+1})\bigr] - f(w_0)
\le -\eta C_1\sum_{t=1}^{T}\frac{1}{\sqrt{t}}
      \EE\bigl[\|\nabla f(w_t)\|^2\bigr]
   + \eta^2 L C_2\sum_{t=1}^{T}\frac{1}{t}.
\tag{14}
\]

Since $f(w_{T+1})\ge f^\star$, rearrange:
\[
\sum_{t=1}^{T}\frac{1}{\sqrt{t}}
      \EE\bigl[\|\nabla f(w_t)\|^2\bigr]
\le
\frac{f(w_0)-f^\star}{\eta C_1}
   + \eta L C_2\sum_{t=1}^{T}\frac{1}{t}.
\tag{15}
\]

\item[Step 15.] \textbf{Lower bound the harmonic sum of $1/\sqrt{t}$.}
Using $\sum_{t=1}^{T}t^{-1/2}\ge 2(\sqrt{T}-1)$,
\[
\frac{1}{T}\sum_{t=1}^{T}
      \EE\bigl[\|\nabla f(w_t)\|^2\bigr]
\le
\frac{1}{T}\frac{1}{2(\sqrt{T}-1)}
\Bigl(
\frac{f(w_0)-f^\star}{\eta C_1}
   + \eta L C_2\,(1+\log T)
\Bigr).
\tag{16}
\]

\item[Step 16.] \textbf{Simplify constants.}
Recall $c_{\min}=1/(\sqrt{G^2/(1-\beta_2)}+\epsilon)$,
so $C_1 = \frac{1}{2}\bigl(\frac{1}{\sqrt{G^2/(1-\beta_2)}+\epsilon}\bigr)$.
Absorbing $\epsilon$ into a constant and using the bound
$\frac{1}{\sqrt{G^2/(1-\beta_2)}+\epsilon}
\le \sqrt{1-\beta_2}/G$,
we obtain
\[
C_1 \ge \frac{\sqrt{1-\beta_2}}{2G},\qquad
C_2 \le \frac{G^2}{(1-\beta_1)^2}\frac{( \sqrt{1-\beta_2}+ \epsilon G)^2}{1-\beta_2}
     \le \frac{G^2}{(1-\beta_1)^2(1-\beta_2)^{1/2}} .
\]

Plugging these bounds into (16) gives the final rate
\[
\frac{1}{T}\sum_{t=1}^{T}
   \EE\bigl[\|\nabla f(w_t)\|^2\bigr]
\le
\underbrace{\frac{2L\bigl(f(w_0)-f^\star\bigr)}{\eta\sqrt{T}}}_{\text{optimization term}}
\;+\;
\underbrace{\frac{2\eta G^2(1+\lambda)}{(1-\beta_1)(1-\beta_2)^{1/2}}
          \frac{\log T}{\sqrt{T}}}_{\text{variance/adaptivity term}} .
\]
This is exactly the statement of Theorem~\ref{thm:adamw}. ∎

\section*{Rate Derivation (With Constants)}
The bound in Theorem~\ref{thm:adamw} can be written succinctly as
\[
\Phi_T \le
\mathcal{O}\!\Bigl(\frac{\log T}{\sqrt{T}}\Bigr),
\]
where the hidden constant depends linearly on
$\displaystyle\frac{G^2}{(1-\beta_1)(1-\beta_2)^{1/2}}$
and on the weight‑decay coefficient $\lambda$.
If $\lambda=0$ the second term reduces to the standard Adam constant.

\section*{Special Cases}
\begin{enumerate}[noitemsep]
    \item \textbf{Pure SGD with weight decay.}
    Setting $\beta_1=\beta_2=0$ gives $m_t=g_t$, $v_t=g_t^2$,
    $\hat m_t=g_t$, $\hat v_t=g_t^2$, and the update reduces to
    $w_{t+1}=w_t-\eta_t\lambda w_t-\eta_t g_t$.
    The bound becomes $O(1/\sqrt{T})$, matching classical SGD analysis.
    \item \textbf{No weight decay ($\lambda=0$).}
    The theorem collapses to the known Adam convergence rate
    $O((\log T)/\sqrt{T})$.
    \item \textbf{Constant learning rate.}
    If $\eta_t\equiv\eta$, the summations in (14)–(15) give a
    $O(1/T)$ bound on the *minimum* gradient norm, but the
    iterates may not converge due to the lack of diminishing steps.
\end{enumerate}

\begin{thebibliography}{9}
\bibitem{Reddi2018}
S.~J. Reddi, S.~Kale, and S.~Kumar, ``On the convergence of adam and
  beyond,'' \emph{International Conference on Learning Representations},
  2018.

\bibitem{Nesterov2004}
Y.~Nesterov, \emph{Introductory Lectures on Convex Optimization:
  A Basic Course}. Springer, 2004.
\end{thebibliography}

\end{document}